% discussion.tex — HapGraph Discussion
% Target: 400+ words, 3+ citations

\hapgraph{} demonstrates that F-statistics and IBD sharing statistics can be
combined in a two-stage integrated pipeline to address a long-standing gap
in population genetic inference: estimating not only the topology and
proportions of an admixture graph but also the timing of each admixture event.
We emphasise that the two stages are deliberately \emph{sequential} rather
than ``jointly'' inferred in the full Bayesian sense: the topology is treated
as fixed when estimating $T$, not as a random variable.
This design choice sacrifices coherent joint uncertainty propagation for
computational tractability at $N \gtrsim 10$ populations.
Here we discuss the design trade-offs that shaped the current implementation,
the limitations revealed by our simulation experiments, and the directions
that would most improve the tool.

\paragraph{Two-stage decoupling: a necessary trade-off.}
The decision to decouple topology inference (Stage I, F-statistics only) from
timing inference (Stage II, IBD) was motivated by preliminary experiments in
which including IBD in the greedy search caused systematic edge over-detection,
consistent with the known confounding between IBD signals and topology ambiguity
\citep{Loh2013}.
This decoupling prevents coherent joint uncertainty propagation — the
topology is treated as fixed when estimating $T$, rather than as a random
variable — and is therefore a limitation of the current design relative to
a fully Bayesian reversible-jump MCMC framework such as
\textsc{AdmixtureBayes} \citep{Nielsen2023}.
However, full RJMCMC over graph space is computationally prohibitive for
$N \gtrsim 10$ populations \citep{Nielsen2023}, whereas \hapgraph{} runs
in seconds to minutes per simulation replicate.
The two-stage approach thus represents a practical compromise between
statistical completeness and scalability, a design philosophy shared by the
wider family of greedy-then-refine methods in phylogenetics.

\paragraph{Edge detection accuracy and the stopping criterion.}
In S2 ($K=1$), the greedy search recovered the exact topology in 17 of 20
seeds (85\%); in the remaining 3 seeds, the search either missed the admixed
target or mislabelled a non-admixed population, each reducing both recall and
precision to 0.85.
The stopping criterion in equation~(\ref{eq:penalty}) is a heuristic that
combines a BIC-inspired penalty with a safety constant; it is not a formally
derived information criterion, because the $N(N-1)/2$ pairwise $F_2$
statistics are correlated across loci, violating the independence assumption
underlying BIC.
The constant 5.0 was chosen empirically to achieve 0\% false positive rate
on pure trees (S1) with the present simulation parameters; it may need
recalibration for datasets with different SNP density or population size.
In S3 ($K=2$, ten populations), the richer F-statistic structure provides
stronger penalty discrimination, and the exact topology was recovered in 19 of
20 seeds (95\%); the single failing seed identified one of the two events
but missed the other, maintaining perfect precision.
We recommend users treat the inferred $K$ as an upper bound in data-sparse
settings and use $F_4$ consistency tests \citep{Patterson2012} to
independently validate detected events.

\paragraph{Alpha uncertainty quantification.}
The 95\% credible interval for $\alpha$ achieved only 10--20\% empirical
coverage across scenarios, substantially below the nominal 95\%.
Two factors compound to produce this severe undercoverage.
First, the $F_3$ method-of-moments delta-method approximation is known to
underestimate uncertainty when $\alpha \approx 0.5$ (where the quadratic
estimator is least sensitive to $F_3$ perturbations) and when the
source populations are imperfectly represented by the NJ tree
(introducing model misspecification not captured by the jackknife
standard error).
Second, our simulation benchmarks used $N_e = 1{,}000$, which produces
elevated genetic drift and inflated pairwise $F_2$ values; under these
conditions the delta-method variance estimate, which assumes approximately
normal $F_3$ statistics, is particularly unreliable.
More principled alpha uncertainty quantification --- for instance, by
bootstrapping the $F_3$ and $F_2$ statistics jointly across blocks, or by
propagating NJ tree uncertainty via a parametric bootstrap --- is a priority
for future development.
We caution users that the reported $\alpha$ confidence intervals should
be treated as approximate guides rather than calibrated 95\% intervals
in the current implementation.

\paragraph{Timing estimation: calibration and scope.}
The 95\% credible interval for $T$ achieved 95--97.5\% empirical coverage,
exceeding the pre-specified 85\% target.
This indicates that the MCMC posterior for timing is somewhat conservative
(over-dispersed), which is preferable to overconfidence.
The IBD timing model uses the corrected expectation
$\mathbb{E}[\bar{L} \mid \bar{L} > L_{\min}] = 50/T + L_{\min}$
(reflecting that IBD between modern individuals decays at rate $1/(2T)$ per
Morgan rather than $1/T$, plus a truncation bias term; $L_{\min} = 2$~cM),
and a noise model that combines sampling variance ($\propto 1/n_k$, where
$n_k$ is the IBD segment count for pair $k$) with a 2~cM background floor.
This formulation is well-specified under the single-pulse assumption and
the NUTS sampler achieves near-perfect $\hat{R}$ in all simulation replicates.
Importantly, the simulation benchmarks used $N_e = 1{,}000$ and 0.5-Morgan
equivalent genome length; under these conditions the mean IBD length is
influenced by a coalescence correction of order $2N_e / n_{\text{haploids}}$
that can bias $T$ estimates toward a narrow range
(see the calibration note in the Methods).
In whole-genome applications such as the 1kGP, where the effective genome
length is $\sim$35~Morgans and $N_e$ is larger, this bias is substantially
reduced.
The 2~cM segment detection threshold of \hapibd{} \citep{Browning2020}
restricts reliable inference to $T \in [5,\, 50]$ generations; admixture
events older than 50 generations produce segments shorter than 2~cM that
\hapibd{} cannot reliably detect, while very recent events ($T < 5$) produce
segments too long to distinguish from background relatedness.
Users analysing deep historical admixture ($T > 50$) should instead
use LD-decay tools such as \alder{} \citep{Loh2013} or
\textsc{GLOBETROTTER} \citep{Hellenthal2014}.

\paragraph{Scope limitations and future work.}
Several limitations of the current study should be noted.
Systematic simulation validation was conducted only for $K \leq 2$; the
accuracy of topology inference and parameter estimation for $K \geq 3$ was
demonstrated qualitatively on the 1kGP real data but not benchmarked against
known ground truth.
Admixture timing was not estimated for the 1kGP application, as it requires
phased IBD pre-computation that is available in principle but was deferred
from the current study.
No direct quantitative comparison to \treemix{} or other baseline tools was
performed; because \treemix{} does not estimate $T$, a comparison of alpha
and topology accuracy would illuminate the cost of the joint framework
relative to the F-statistics-only baseline, and is a priority for future work.
Finally, the current model assumes a single admixture pulse per target
population; extending the IBD likelihood to accommodate multiple pulses
or continuous gene flow would address a known limitation of pulse models
in populations with prolonged admixture histories \citep{Loh2013}.
